\scanpage{169}
% archive foliation: 168
\margin{(Diviseurs de \uline{fonctions}}
Soit $X$ une variété irréductible \uline{normale} déf/$k$

$K=k(X)$ le corps des fonctions rationnelles/$k$ sur $X$

Soit $V$ un \deleted{\unsure{hypersurface irréductible}} \inserted{diviseur premier déf.} sur $k$ sur $X$

Alors l'ensemble des $f\in K$ \uline{définies} en $V$ forment un anneau de valuation discrète, \deleted{$\unsure{O_V}$} $\underline O_{(V)}$.

En effet, c'est une \deleted{\illegible{1}} localité normale, dont la dimension de Krull est 1.

Soit $v_V$ la valuation \inserted{de $K$} associée à cet anneau de valuation de $K$. Appelons \uline{diviseur} \inserted{déf/$k$} sur $X$ toute comb. lin. formelle des diviseurs premiers sur \unsure{déf}/$k$. Posons, pour tout $f\in K$
\[
\operatorname{div}f=\sum_{\substack{V\,\mathrm{div.}\\\mathrm{pr}/k}}v_V(f).V
\]
Cette somme a un sens, car pour $f$ donné, il n'y a qu'un nb fini de termes non nuls; \deleted{\unsure{ceci étant}} pour le voir, on est ramené au cas \deleted{\unsure{normal}} affine, alors les diviseurs pr. déf/$k$ correspondent aux idéaux premiers non nuls minimaux de l'anneau normal noethérien $A=k[X]$, et la chose est bien connue.

\uline{Théorème} \deleted{$\unsure{f\in K}$ \illegible{3} \unsure{partout}} \inserted{$\operatorname{div}f\geq0.\Longleftrightarrow f$ est partout définie}

Car. résulte de la théorie des anneaux normaux;
\[
A=\bigcap A_{\mathfrak p}\qquad(\mathfrak p\text{ premiers minimaux})\quad\text{si }A\text{ est normal.}
\]
En effet, on est aussitôt ramené au cas affine, et $\operatorname{div}f\geq0\Longleftrightarrow f\in\bigcap A_{\mathfrak p}$.

